\section{Economic Environment and Stand-Alone Exposures}\label{sec:model}

Begin with the familiar finite-dimensional factor model, in which an exposure
is a vector of $K$ factor loadings in $\R^K$.
Firm $i$ belongs to a finite universe of $n$ firms and has a population
\emph{characteristic law} $C_i$ on an embedding space
$(\Omega,d_\Omega)$, which summarizes the distribution of its information.
In the application, observed articles produce the empirical law
$\widehat C_i$ as a proxy for $C_i$.
Think of the firm's \emph{stand-alone exposure} $\xi_i$ as the factor-loading
vector implied by its own information before any peer adjustment.
Peer adjustment maps $\xi_i$ into the latent, model-implied
\emph{peer-adjusted exposure} $B_i$, and a maintained factor bridge then links
$B_i$ to observed centered excess returns.
The economic sequence is therefore characteristic law, stand-alone exposure,
peer-adjusted exposure, and return.

The firm's economic object is an exposure, not a return.
The quadratic criterion introduced in the next section represents, in reduced
form, the costs of operational, financing, or portfolio reconfiguration.
It is an as-if adjustment problem and does not require firms literally to choose
factor loadings each period.
Keeping $\xi_i$ and $B_i$ separate lets the model ask how much of the exposure
that enters returns reflects the firm's own information and how much reflects
alignment with other firms.

Observed information is not itself an exposure: embedding locations describe
an information distribution, whereas factor loadings measure sensitivity to
common shocks.
A transmission map is therefore needed to connect the distribution-valued
characteristic law to stand-alone exposure in factor-loading units.
Formally, draw $X_i\sim C_i$ and let $U_i$ collect idiosyncratic transmission
randomness.
A common measurable map $T$ converts the information draw and transmission shock
into the exposure space:

\begin{equation}\label{eq:primitive-exposure}
  \xi_i = T(X_i,U_i).
\end{equation}

Together, $C_i$, $U_i$, and $T$ generate the stand-alone exposure $\xi_i$.
The characteristic law is measured in embedding-space units, whereas $\xi_i$
is measured in factor-exposure units.
The map $T$, its latent inputs, and the cross-firm coupling of the resulting
stand-alone exposures are not identified from the return panel.
No peer response has yet entered Equation~\eqref{eq:primitive-exposure}.

Moving from $\xi_i$ to $B_i$ requires a peer average and a relative adjustment
weight.
Let $W$ be the peer matrix that describes how each firm weights the other firms,
and let
$\lambda\geq0$ be the dimensionless weight on peer alignment relative to the
unit cost of departing from $\xi_i$.
Because peer alignment is averaging over an economic neighborhood rather than
forming a signed contrast, each row must use nonnegative weights, exclude the
firm itself, and sum to one.

\begin{definition}[Row-stochastic, zero-diagonal interaction
  matrix]\label{def:interaction-matrix}
  An interaction matrix is any $W \in \R^{n \times n}$ with $W_{ij} \geq 0$ for
  all $i$ and $j$, $W_{ii} = 0$ for all $i$, and $\sum_j W_{ij} = 1$ for all
  $i$.
  A firm never interacts with itself, weights every other firm nonnegatively,
  and spends exactly one unit of interaction weight on the remaining firms.
\end{definition}

The definition makes $\sum_{j} W_{ij} B_j$ a peer-weighted average in the same
factor-exposure units as $B_i$.
It restricts the economic role of each row but does not select its weights.
Later, Wasserstein geometry will align the empirical characteristic laws, and
simplex weights will form a target-specific barycentric interaction field.
That field will be represented by a matrix satisfying the definition above;
\Cref{sec:interaction-operator} supplies the formal construction.
For now, the economic environment takes $W$ as a fixed admissible peer matrix.

In the finite-dimensional case, $\xi_i$ and $B_i$ are ordinary vectors of $K$
factor loadings, and peer adjustment operates coordinate by coordinate.
To cover either finitely or countably many risk directions in one statement, we
now let the exposures take values in a real separable Hilbert space $\mathcal H$
that generalizes $\R^K$.
Let $F_t\in\mathcal H$ be a centered common factor innovation with covariance
operator $\Gamma$, and let $e_t^i$ be a centered idiosyncratic return component.
The maintained return bridge specifies how the peer-adjusted exposure enters
centered excess returns:

\begin{equation}\label{eq:hilbert-return}
  \widetilde r_t^i = \inner{B_i,F_t}_{\mathcal H}+e_t^i.
\end{equation}

Here $B_i$ has loading units, $F_t$ has factor-innovation units, and their inner
product has return units.
This bridge requires the maintained conditions that $e_t^i$ is
square-integrable, orthogonal to $F_t$, and has zero cross-firm covariance.

The environment now contains the stand-alone exposure $\xi$, the peer-adjusted
exposure $B$, an admissible peer matrix $W$, and the relative adjustment weight
$\lambda$.
The equilibrium is solved pointwise for each realization of $\xi$.
Consequently, $B$ is random whenever $\xi$ is random, even when $W$ and
$\lambda$ are fixed.
The next section asks whether a transparent firm-level objective maps
stand-alone exposures into a unique peer-adjustment equilibrium.
